\vspace{-2pt}
\section{Conclusion}
\vspace{-3pt}

In this paper, we present CEIL, a novel and general Imitation Learning framework applicable to a wide range of IL settings, including \textit{C/S-on/off-LfD/LfO} and few-shot IL settings. 
This is achieved by explicitly decoupling the imitation policy into 1) a contextual policy, learned with the self-supervised hindsight information matching objective, and 2) a latent variable, inferred by performing the IL expert matching objective. 
Compared to prior baselines, our results show that CEIL is more sample-efficient in most of the online IL tasks and achieves better or competitive performances in offline~tasks.  





\textbf{Limitations and future work.} 
Our primary aim behind this work is to develop a simple and scalable IL method. We believe that CEIL makes an important step in that direction. Admittedly, we also find some limitations of CEIL: 1) Offline results generally outperform online results, especially in the LfO setting. The main reason is that CEIL lacks explicit exploration bounds, thus future work could explore the exploration ability of online CEIL. 
2) The trajectory self-consistency cannot be applied to cross-embodiment agents once the two embodiments/domains have different state spaces or action spaces. 
Considering such a cross-embodiment setting, a typical approach is to serialize state/action from different modalities into a flat sequence of tokens. We also remark that CEIL is compatible with such a tokenization approach, and thus suitable for IL tasks with different action/state spaces. %, yet it is out of the scope of the current work.
Thus, we encourage the future exploration of generalized IL methods across different embodiments.

% to explore cross-embodiment IL in future. 




% Our CEIL offers the potential to scale the IL method to general and broadly applicable

% The key contribution of this work is a simple and scalable method 


% In future work, we would like to give attention to issues in situations where the data quality could not be lower bounded. Although learning objectives expressed and optimized in the supervised ..., our results elicit further in-depth reflection on 